%!TEX root = ../thesis.tex

\chapter{Conclusion}

\section{Aim and Evidence Base}
% <<SECTION_1_BEGIN>>
The central question was how precision and hardware/backend choices affect a Riemann-solver-based HRSC method on a controlled Euler validation suite. Evidence came from one- and two-dimensional Euler cases, matched CPU/GPU saved-output comparisons, compiler and branch-rule variations, direct fp32/fp64 reference scaling, and region-aware precision diagnostics.
% <<SECTION_1_END>>

\section{Key Findings}
% <<SECTION_2_BEGIN>>
Matched strict-IEEE HLLC CPU/GPU runs showed no observed saved-state divergence within each tested case, with zero \(L_1\), \(L_\infty\), and \(\mathrm{ULP}_{\max}\). The CSC Linux strict-IEEE rerun covers Toro3 and Toro5 final state plus four checkpoints; matched saved-output coverage for Sod, LW3, and LW12 uses local Linux/WSL paths. This match holds for saved outputs and checkpoints, not every intermediate value inside a time step. Implementation sensitivity was axis-dependent: O2/O3 comparisons were bit-identical where tested, fast-math changed non-stationary saved states, and HLLC--Rusanov was deliberate method variation. In these saved final states, direct fp32/fp64 gaps stayed below the available reference/discretisation scale; for LW12 at \(400^2\), the density ratio was \(1.30\times10^{-4}\) and the pressure ratio was \(1.13\times10^{-4}\). Verificarlo \texttt{p32}, a virtual rather than IEEE fp32 mode, then located remaining precision pressure on LW3 fronts. Together these findings give a bounded Euler baseline. The later precision and hardware study builds on it.
% <<SECTION_2_END>>

\section{Limitation and Next Step}
% <<SECTION_3_BEGIN>>
Detailed limitations are listed in Section~\ref{sec:ch6-limits}. Report~2 will extend the drift-time-series framework to ideal-MHD test cases with divergence control.
% <<SECTION_3_END>>
